\documentclass[10pt,a4paper]{amsart}
\usepackage{amsmath,amssymb,amsthm,graphicx}
\usepackage{verbatim}
\newenvironment{reference}{\comment}{\endcomment}
%\newenvironment{reference}{}{}
\newenvironment{slogan}{\comment}{\endcomment}
\newenvironment{history}{\comment}{\endcomment}
% For commutative diagrams we use Xy-pic
\usepackage[all]{xy}

% We use 2cell for 2-commutative diagrams.
\xyoption{2cell}
% Theorem environments.
%
\theoremstyle{plain}
\newtheorem{theorem}[subsection]{Theorem}
\newtheorem{proposition}[subsection]{Proposition}
\newtheorem{lemma}[subsection]{Lemma}

\theoremstyle{definition}
\newtheorem{definition}[subsection]{Definition}
\newtheorem{example}[subsection]{Example}
\newtheorem{exercise}[subsection]{Exercise}
\newtheorem{situation}[subsection]{Situation}

\theoremstyle{remark}
\newtheorem{remark}[subsection]{Remark}
\newtheorem{remarks}[subsection]{Remarks}

\numberwithin{equation}{subsection}

% Macros
%
\def\lim{\mathop{\mathrm{lim}}\nolimits}
\def\colim{\mathop{\mathrm{colim}}\nolimits}
\def\Spec{\mathop{\mathrm{Spec}}}
\def\Hom{\mathop{\mathrm{Hom}}\nolimits}
\def\Ext{\mathop{\mathrm{Ext}}\nolimits}
\def\SheafHom{\mathop{\mathcal{H}\!\mathit{om}}\nolimits}
\def\SheafExt{\mathop{\mathcal{E}\!\mathit{xt}}\nolimits}
\def\Sch{\mathit{Sch}}
\def\Mor{\mathop{\mathrm{Mor}}\nolimits}
\def\Ob{\mathop{\mathrm{Ob}}\nolimits}
\def\Sh{\mathop{\mathit{Sh}}\nolimits}
\def\NL{\mathop{N\!L}\nolimits}
\def\CH{\mathop{\mathrm{CH}}\nolimits}
\def\proetale{{pro\text{-}\acute{e}tale}}
\def\etale{{\acute{e}tale}}
\def\QCoh{\mathit{QCoh}}
\def\Ker{\mathop{\mathrm{Ker}}}
\def\Im{\mathop{\mathrm{Im}}}
\def\Coker{\mathop{\mathrm{Coker}}}
\def\Coim{\mathop{\mathrm{Coim}}}

% Boxtimes
%
\DeclareMathSymbol{\boxtimes}{\mathbin}{AMSa}{"02}

%
% Macros for moduli stacks/spaces
%
\def\QCohstack{\mathcal{QC}\!\mathit{oh}}
\def\Cohstack{\mathcal{C}\!\mathit{oh}}
\def\Spacesstack{\mathcal{S}\!\mathit{paces}}
\def\Quotfunctor{\mathrm{Quot}}
\def\Hilbfunctor{\mathrm{Hilb}}
\def\Curvesstack{\mathcal{C}\!\mathit{urves}}
\def\Polarizedstack{\mathcal{P}\!\mathit{olarized}}
\def\Complexesstack{\mathcal{C}\!\mathit{omplexes}}
% \Pic is the operator that assigns to X its picard group, usage \Pic(X)
% \Picardstack_{X/B} denotes the Picard stack of X over B
% \Picardfunctor_{X/B} denotes the Picard functor of X over B
\def\Pic{\mathop{\mathrm{Pic}}\nolimits}
\def\Picardstack{\mathcal{P}\!\mathit{ic}}
\def\Picardfunctor{\mathrm{Pic}}
\def\Deformationcategory{\mathcal{D}\!\mathit{ef}}
\usepackage{cleveref}
\crefname{originallemma}{original lemma}{original lemmas}
\Crefname{originallemma}{Original lemma}{Original lemmas}
\makeatletter
\newlabel{lemma-base-change-universally-closed}{{lemma-base-change-universally-closed}{source}{}{sourcepointer-1}{}}
\newlabel{lemma-base-change-universally-closed@cref}{{[originallemma][0][]lemma-base-change-universally-closed}{[1][0][]source}{}{}{}}
\makeatother
\crefname{originallemma}{original lemma}{original lemmas}
\Crefname{originallemma}{Original lemma}{Original lemmas}
\makeatletter
\newlabel{lemma-compose-after-separated}{{lemma-compose-after-separated}{source}{}{sourcepointer-2}{}}
\newlabel{lemma-compose-after-separated@cref}{{[originallemma][0][]lemma-compose-after-separated}{[1][0][]source}{}{}{}}
\makeatother
\crefname{originallemma}{original lemma}{original lemmas}
\Crefname{originallemma}{Original lemma}{Original lemmas}
\makeatletter
\newlabel{lemma-composition-universally-closed}{{lemma-composition-universally-closed}{source}{}{sourcepointer-3}{}}
\newlabel{lemma-composition-universally-closed@cref}{{[originallemma][0][]lemma-composition-universally-closed}{[1][0][]source}{}{}{}}
\makeatother
\crefname{originallemma}{original lemma}{original lemmas}
\Crefname{originallemma}{Original lemma}{Original lemmas}
\makeatletter
\newlabel{lemma-finite-type-permanence}{{lemma-finite-type-permanence}{source}{}{sourcepointer-4}{}}
\newlabel{lemma-finite-type-permanence@cref}{{[originallemma][0][]lemma-finite-type-permanence}{[1][0][]source}{}{}{}}
\makeatother
\crefname{originallemma}{original lemma}{original lemmas}
\Crefname{originallemma}{Original lemma}{Original lemmas}
\makeatletter
\newlabel{lemma-quasi-compact-permanence}{{lemma-quasi-compact-permanence}{source}{}{sourcepointer-5}{}}
\newlabel{lemma-quasi-compact-permanence@cref}{{[originallemma][0][]lemma-quasi-compact-permanence}{[1][0][]source}{}{}{}}
\makeatother
\crefname{originallemma}{original lemma}{original lemmas}
\Crefname{originallemma}{Original lemma}{Original lemmas}
\makeatletter
\newlabel{lemma-semi-diagonal}{{lemma-semi-diagonal}{source}{}{sourcepointer-6}{}}
\newlabel{lemma-semi-diagonal@cref}{{[originallemma][0][]lemma-semi-diagonal}{[1][0][]source}{}{}{}}
\makeatother
\title{Universal closedness and properness are preserved when the target is separated}
\author{The Stacks Project Authors}
\date{stacks-project:stacks-morphisms.tex:local-snapshot-20261012; GNU FDL 1.2 or later}
\begin{document}
\maketitle
\noindent\textit{Editorial note.} All delivered mathematical text is the author's own, copied from the pinned Stacks Project chapter source sources/stacks-algebraic-stacks-morphisms/pinned-original-source.tex; no statement or proof step was rewritten, reordered or removed.\par\medskip
\noindent\textit{Editorial note.} Preamble provenance: the Stacks Project keeps one single shared preamble for every chapter; this card takes it directly from the staged sources/stacks-algebraic-stacks-morphisms/preamble.tex as three verbatim spans (lines 8-12, 14-18 and 154-217), so no derived file was written and nothing was added or redefined. The wrapper class and the base packages are supplied by the assembler.\par\medskip
\noindent\textit{Editorial note.} Bibliography provenance: neither the counted statement nor its proof contains a citation (this lane rescans the whole span before writing the card), so the card carries no bibliography and no entry was reconstructed.\par\medskip
\noindent\textit{Editorial note.} Selection provenance: this card was opened under the file-based control contract (CONTROL.md), with the mechanical pre-card gate run on its statement and proof interval before the card was written.\par\medskip
\noindent\textit{Allowed standard prerequisites (editorial).}\begin{itemize}
\item The statement that the graph of a morphism into a separated stack is proper. Invoked inside the counted proof (source lines 8677-8678). Makes the first factor of the factorization universally closed.
\item The base change permanence of universal closedness. Invoked inside the counted proof (source lines 8679-8682). Makes the projection universally closed.
\item The composition permanence of universal closedness. Invoked inside the counted proof (source lines 8683-8685). Combines the two factors.
\item The separatedness and finite type permanence statements. Invoked at the end of the counted proof (source lines 8686-8690). Upgrades the conclusion from universally closed to proper.
\end{itemize}\medskip
\ifdefined\hypertarget\hypertarget{sourcepointer-1}{}\fi
\noindent\textit{Source reference.} Original lemma lemma-base-change-universally-closed: sources/stacks-algebraic-stacks-morphisms/pinned-original-source.tex, source lines 2726--2726. Navigation pointer only: the counted text refers to the Stacks item lemma-base-change-universally-closed, which is an existing labeled result of this same chapter and is not part of this card. The pointer carries the target's own label as its number and its definition line as its span; its content is not reproduced here.\par\medskip
\ifdefined\hypertarget\hypertarget{sourcepointer-2}{}\fi
\noindent\textit{Source reference.} Original lemma lemma-compose-after-separated: sources/stacks-algebraic-stacks-morphisms/pinned-original-source.tex, source lines 929--929. Navigation pointer only: the counted text refers to the Stacks item lemma-compose-after-separated, which is an existing labeled result of this same chapter and is not part of this card. The pointer carries the target's own label as its number and its definition line as its span; its content is not reproduced here.\par\medskip
\ifdefined\hypertarget\hypertarget{sourcepointer-3}{}\fi
\noindent\textit{Source reference.} Original lemma lemma-composition-universally-closed: sources/stacks-algebraic-stacks-morphisms/pinned-original-source.tex, source lines 2736--2736. Navigation pointer only: the counted text refers to the Stacks item lemma-composition-universally-closed, which is an existing labeled result of this same chapter and is not part of this card. The pointer carries the target's own label as its number and its definition line as its span; its content is not reproduced here.\par\medskip
\ifdefined\hypertarget\hypertarget{sourcepointer-4}{}\fi
\noindent\textit{Source reference.} Original lemma lemma-finite-type-permanence: sources/stacks-algebraic-stacks-morphisms/pinned-original-source.tex, source lines 3608--3608. Navigation pointer only: the counted text refers to the Stacks item lemma-finite-type-permanence, which is an existing labeled result of this same chapter and is not part of this card. The pointer carries the target's own label as its number and its definition line as its span; its content is not reproduced here.\par\medskip
\ifdefined\hypertarget\hypertarget{sourcepointer-5}{}\fi
\noindent\textit{Source reference.} Original lemma lemma-quasi-compact-permanence: sources/stacks-algebraic-stacks-morphisms/pinned-original-source.tex, source lines 2044--2044. Navigation pointer only: the counted text refers to the Stacks item lemma-quasi-compact-permanence, which is an existing labeled result of this same chapter and is not part of this card. The pointer carries the target's own label as its number and its definition line as its span; its content is not reproduced here.\par\medskip
\ifdefined\hypertarget\hypertarget{sourcepointer-6}{}\fi
\noindent\textit{Source reference.} Original lemma lemma-semi-diagonal: sources/stacks-algebraic-stacks-morphisms/pinned-original-source.tex, source lines 808--808. Navigation pointer only: the counted text refers to the Stacks item lemma-semi-diagonal, which is an existing labeled result of this same chapter and is not part of this card. The pointer carries the target's own label as its number and its definition line as its span; its content is not reproduced here.\par\medskip
\begin{lemma}
\label{lemma-universally-closed-permanence}
Consider a commutative diagram
$$
\xymatrix{
\mathcal{X} \ar[rr] \ar[rd] & &
\mathcal{Y} \ar[ld] \\
& \mathcal{Z} &
}
$$
of algebraic stacks.
\begin{enumerate}
\item If $\mathcal{X} \to \mathcal{Z}$ is universally closed and
$\mathcal{Y} \to \mathcal{Z}$ is separated,
then the morphism $\mathcal{X} \to \mathcal{Y}$ is universally closed.
In particular, the image of $|\mathcal{X}|$ in $|\mathcal{Y}|$ is closed.
\item If $\mathcal{X} \to \mathcal{Z}$ is proper and
$\mathcal{Y} \to \mathcal{Z}$ is separated, then
the morphism $\mathcal{X} \to \mathcal{Y}$ is proper.
\end{enumerate}
\end{lemma}

\begin{proof}
Assume $\mathcal{X} \to \mathcal{Z}$ is universally closed and
$\mathcal{Y} \to \mathcal{Z}$ is separated.
We factor the morphism as
$\mathcal{X} \to \mathcal{X} \times_\mathcal{Z} \mathcal{Y} \to \mathcal{Y}$.
The first morphism is proper (Lemma \ref{lemma-semi-diagonal})
hence universally closed.
The projection $\mathcal{X} \times_\mathcal{Z} \mathcal{Y} \to \mathcal{Y}$
is the base change of a universally closed morphism and hence
universally closed, see
Lemma \ref{lemma-base-change-universally-closed}.
Thus $\mathcal{X} \to \mathcal{Y}$ is universally closed as the composition
of universally closed morphisms, see
Lemma \ref{lemma-composition-universally-closed}.
This proves (1). To deduce (2) combine (1) with
Lemmas \ref{lemma-compose-after-separated},
\ref{lemma-quasi-compact-permanence}, and
\ref{lemma-finite-type-permanence}.
\end{proof}

\end{document}
